\documentclass[a4paper,12pt]{article}

% 页面设置
\usepackage[top=2.54cm, bottom=2.54cm, left=1.91cm, right=1.91cm]{geometry}

% 中文支持
\usepackage[UTF8]{ctex}
\usepackage{fontspec}

% 数学公式支持
\usepackage{amsmath}
\usepackage{amssymb}
\usepackage{amsfonts}
\usepackage{amsthm}

\usepackage{enumitem}

% 定理环境定义
% \newtheorem*{thm}{Thm}
% \newtheorem*{define}{Def}
% \newtheorem{lemma}{引理}
% \newtheorem*{corollary}{Corollary}
% \newtheorem*{example}{Eg}
% \newtheorem*{rmk}{Rmk}
% \newtheorem*{prop}{Prop}
\newtheorem*{证}{证}
\newtheorem*{解}{解}

% 图形支持
\usepackage{graphicx}
\usepackage{tikz}

% 超链接支持
\usepackage{hyperref}
\hypersetup{
    colorlinks=true,
    linkcolor=blue,
    filecolor=magenta,
    urlcolor=cyan,
}

% 行距设置
\linespread{1.25}

% 标题信息
\title{Week 1 Homework (March 5)}
\author{袁子强}
% \date{}

\begin{document}

\maketitle

\section*{Section 8.1}

13. 计算以向量 $a=p-3q+r,\,b=2p+q-3r$ 和 $c=p+2q+r$ 为棱的平行六面体的体积，这里 $p,q$ 和 $r$ 是互相垂直的单位向量。

\begin{解}
$$\text{Vol}(a,b,c)=\left<a,b\times c\right>=\begin{vmatrix}
        1 & -3 & 1  \\
        2 & 1  & -3 \\
        1 & 2  & 1
    \end{vmatrix}=1+9+4-1+6+6=25$$
\end{解}

17. 已知 $\overrightarrow{OA} = 2\boldsymbol{i} + 4\boldsymbol{j} + \boldsymbol{k},\,\overrightarrow{OB} = 3\boldsymbol{i} + 7\boldsymbol{j} + 5\boldsymbol{k},\,\overrightarrow{OC} = 4\boldsymbol{i} + 10\boldsymbol{j} + 9\boldsymbol{k}$，问 $A, B, C$ 三点是否共线？

\begin{解}
即判断 $\begin{cases}\overrightarrow{AB}=\overrightarrow{OB}-\overrightarrow{OA}=\boldsymbol{i} + 3\boldsymbol{j} + 4\boldsymbol{k}\\\overrightarrow{AC}=\overrightarrow{OC}-\overrightarrow{OA}=2\boldsymbol{i} + 6\boldsymbol{j} + 8\boldsymbol{k}\end{cases}$ 是否共线。

显然，$\overrightarrow{AC}=2\overrightarrow{AB}$

即 $\overrightarrow{AB}, \overrightarrow{AC}$ 共线。

即 $A, B, C$ 三点共线。
\end{解}

24. 计算顶点为 $A(2, -1, 1),\,B(5, 5, 4),\,C(3, 2, -1),\,D(4, 1, 3)$ 的四面体的体积。

\begin{解}
即求由 $\begin{cases}
        \overrightarrow{AB}=3\boldsymbol{i} + 6\boldsymbol{j} + 3\boldsymbol{k} \\
        \overrightarrow{AC}=1\boldsymbol{i} + 3\boldsymbol{j} - 2\boldsymbol{k} \\
        \overrightarrow{AD}=2\boldsymbol{i} + 2\boldsymbol{j} + 2\boldsymbol{k}
    \end{cases}$ 为棱的平行六面体的体积的 $\frac{1}{6}$。

$$\text{Vol}(\overrightarrow{AB},\overrightarrow{AC},\overrightarrow{AD})=\left|\left<\overrightarrow{AB},\overrightarrow{AC}\times \overrightarrow{AD}\right>\right|=\left|\begin{vmatrix}
        3 & 6 & 3  \\
        1 & 3 & -2 \\
        2 & 2 & 2
    \end{vmatrix}\right|=\left|18-24+6-18-12+12\right|=18$$

即四面体的体积为 $3$。
\end{解}

28. 求点 $P(4, -3, 5)$ 到坐标原点以及到各坐标轴的距离。

\begin{解}
到原点距离：$|OP|=\sqrt{4^2+(-3)^2+5^2}=\sqrt{50}=5\sqrt{2}$

$x$ 轴：$\begin{cases}y=0\\z=0\end{cases}=t(1,0,0),\,y$ 轴：$\begin{cases}x=0\\z=0\end{cases}=t(0,1,0),\,z$ 轴：$\begin{cases}x=0\\y=0\end{cases}=t(0,0,1)$

所以 $\begin{cases}
        \text{d}(P,x)=\frac{\left|\overrightarrow{OP}\times(1,0,0)\right|}{\left|(1,0,0)\right|}=\begin{vmatrix}
                                                                                                     i & j  & k \\
                                                                                                     4 & -3 & 5 \\
                                                                                                     1 & 0  & 0
                                                                                                 \end{vmatrix}=|5j+3k|=\sqrt{34} \\
        \text{d}(P,y)=\frac{\left|\overrightarrow{OP}\times(0,1,0)\right|}{\left|(0,1,0)\right|}=\begin{vmatrix}
                                                                                                     i & j  & k \\
                                                                                                     4 & -3 & 5 \\
                                                                                                     0 & 1  & 0
                                                                                                 \end{vmatrix}=|4k-5i|=\sqrt{41} \\
        \text{d}(P,z)=\frac{\left|\overrightarrow{OP}\times(0,0,1)\right|}{\left|(0,0,1)\right|}=\begin{vmatrix}
                                                                                                     i & j  & k \\
                                                                                                     4 & -3 & 5 \\
                                                                                                     0 & 0  & 1
                                                                                                 \end{vmatrix}=|-3i-4j|=5
    \end{cases}$
\end{解}

29. 在 $Oyz$ 平面上求一点 $P$，使它与三已知点 $A(3, 1, 2),\,B(4, -2, -2)$ 及 $C(0, 5, 1)$ 等距离。

\begin{解}
$Oyz=\{x=0\}$，设 $P(0,y,z)$，则 $\begin{cases}
        \overrightarrow{PA}=3i,(1-y)j,(2-z)k   \\
        \overrightarrow{PB}=4i,(-2-y)j,(-2-z)k \\
        \overrightarrow{PC}=0i,(5-y)j,(1-z)k
    \end{cases}$。

由题意知，$|PA|=|PB|=|PC|$，即$|PA|^2=|PB|^2=|PC|^2$。

即 $9+1-2y+y^2+4-4z+z^2=16+4+4y+y^2+4+4z+z^2=25-10y+y^2+1-2z+z^2$

即 $\begin{cases}
        3y+4z=-5 \\
        7y+3z=1
    \end{cases}$，解得 $\begin{cases}
        y=1 \\
        z=-2
    \end{cases}$，即 $P(0,1,-2)$。
\end{解}

\section*{Section 8.2}

12. 求两平面 $2x - y + z - 7 = 0$ 和 $x + y + 2z - 11 = 0$ 所成两个二面角的平分面。

\begin{解}
即求到二面距离相等的点集。设 $P(x,y,z)$，即求 $\frac{|2x-y+z-7|}{\sqrt{4+1+1}}=\frac{|x+y+2z-11|}{\sqrt{1+1+4}}$。

即 $2x-y+z-7=x+y+2z-11$ 或 $2x-y+z-7+x+y+2z-11=0$

解得 $x-2y-z+4=0$ 或 $x+z-6=0$
\end{解}

13. 在平面 $x + y + z - 1 = 0$ 与三坐标平面所围成的四面体内求一点，使它到四个面的距离相等。

\begin{解}
即求到四面距离相等的点，设为 $P(x,y,z)$，且满足 $x + y + z - 1 < 0$。

即 $x=y=z=\frac{|x+y+z-1|}{\sqrt{3}}$

即 $x=\left|\sqrt{3}x-\frac{1}{\sqrt{3}}\right|$，且 $0<x<\frac{1}{3}$

即 $x=\frac{1}{\sqrt{3}}-\sqrt{3}x$

解得 $x=\frac{1}{3+\sqrt{3}}=\frac{3-\sqrt{3}}{6}$

得 $P\left(\frac{3-\sqrt{3}}{6},\frac{3-\sqrt{3}}{6},\frac{3-\sqrt{3}}{6}\right)$
\end{解}

15. 分别求出满足下列条件的直线方程：\\
(1) 过点 $(0, 2, 4)$ 而与两平面 $x + 2z = 1$, $y - 3z = 2$ 平行；\\
(2) 过点 $(-1, 2, 1)$ 且平行于直线
$$\begin{cases}
        x + y - 2z - 1 = 0, \\
        x + 2y - z + 1 = 0;
    \end{cases}$$

\begin{解}
(1) $\overrightarrow{r_0}=(0,2,4),\,\overrightarrow{n_1}=(1,0,2),\,\overrightarrow{n_2}=(0,1,-3)$

则 $\overrightarrow{n_1}\times\overrightarrow{n_2}=\begin{vmatrix}
        i & j & k  \\
        1 & 0 & 2  \\
        0 & 1 & -3
    \end{vmatrix}=-2i+3j+k$

得直线方程 $\frac{x}{-2}=\frac{y-2}{3}=z-4$

(2) $\overrightarrow{r_0}=(-1,2,1),\,\overrightarrow{n_1}=(1,1,-2),\,\overrightarrow{n_2}=(1,2,-1)$

则 $\overrightarrow{n_1}\times\overrightarrow{n_2}=\begin{vmatrix}
        i & j & k  \\
        1 & 1 & -2 \\
        1 & 2 & -1
    \end{vmatrix}=-1+4i-2j+j+2k-k=3i-j+k$

得直线方程 $\frac{x+1}{3}=2-y=z-1$
\end{解}

19. 证明下列各组直线互相平行，并求它们间的距离。\\
(1) $\frac{x + 2}{-2} = \frac{y - 1}{3} = z$ 和 $\begin{cases} x + y - z = 0, \\ x - y - 5z - 8 = 0; \end{cases}$\\
(2) $\frac{x + 7}{3} = \frac{y - 5}{-1} = \frac{z - 9}{4}$ 和 $\begin{cases} 2x + 2y - z - 10 = 0, \\ x - y - z - 22 = 0. \end{cases}$

\begin{解}
(1) 方向向量 $
    v_1=(-2,3,1),\,v_2=\begin{vmatrix}
        i & j  & k  \\
        1 & 1  & -1 \\
        1 & -1 & -5
    \end{vmatrix}=-6i+4j-2k=2(-3,2,-1)
$，与 $v_1$ 不成比例，即两直线不平行。

$$???$$

取两直线上的点 取两直线上的点 $R_1=(-2,1,0),\,R_2=(1,-2,-1)$，则 $\overrightarrow{R_1R_2}=(3,-3,-1)$

所以 $d=\left|\left<\overrightarrow{R_1R_2},\frac{v_1\times v_2}{\left|v_1\times v_2\right|}\right>\right|$

其中 $v_1\times v_2=\begin{vmatrix}
        i  & j & k  \\
        -2 & 3 & 1  \\
        -3 & 2 & -1
    \end{vmatrix}=-5i-5j+5k$

所以 $d=\left|\left<\overrightarrow{R_1R_2},\left(-\frac{\sqrt{3}}{3},-\frac{\sqrt{3}}{3},\frac{\sqrt{3}}{3}\right)\right>\right|=\frac{\sqrt{3}}{3}$

(2) 方向向量 $\begin{cases}
        v_1=(3,-1,4) \\
        v_2=\begin{vmatrix}
                i & j  & k  \\
                2 & 2  & -1 \\
                1 & -1 & -1
            \end{vmatrix}=-3i+j-4k=-v_1
    \end{cases}$，即两直线平行。

取两直线上的点 $R_1=(-7,5,9),\,R_2=(0,-4,-18)$，则 $\overrightarrow{R_1R_2}=(7,-9,-27)$

则间距为 $\frac{\left|\overrightarrow{R_1R_2}\times(3,-1,4)\right|}{\sqrt{9+1+16}}=\frac{\left|\begin{vmatrix}
            i & j  & k   \\
            7 & -9 & -27 \\
            3 & -1 & 4
        \end{vmatrix}\right|}{\sqrt{26}}=\frac{\sqrt{(-63)^2+(-109)^2+20^2}}{\sqrt{26}}=\frac{\sqrt{16250}}{\sqrt{26}}=25$
\end{解}

27. 求点 $(-1, 2, 0)$ 在平面 $x + 2y - z + 1 = 0$ 上的投影。

\begin{解}
平面法向量：$v=(1,2,-1)$，则即求直线 $L:\,\overrightarrow{r}=(-1,2,0)+t(1,2,-1)$ 与平面的交点。

即 $-1+t+4+4t+t+1=0$，解得 $t=-\frac{2}{3}$。

即点 $(-1, 2, 0)$ 在平面 $x + 2y - z + 1 = 0$ 上的投影为 $(-\frac{5}{3},\frac{2}{3},\frac{2}{3})$
\end{解}

31. 试求 $\lambda$，使两直线 $\frac{x - 1}{\lambda} = \frac{y + 4}{5} = \frac{z - 3}{-3}$ 与 $\frac{x + 3}{3} = \frac{y - 9}{-4} = \frac{z + 14}{7}$ 相交，并求出交点和由此两直线确定的平面方程。

\begin{解}
参数方程：$L_1\begin{cases}
        x=1+\lambda t \\
        y=-4+5t       \\
        z=3-3t
    \end{cases},\,L_2\begin{cases}
        x=-3+3s \\
        y=9-4s  \\
        z=-14+7s
    \end{cases}$。

联立，得：$\begin{cases}
        1+\lambda t=-3+3s \\
        -4+5t=9-4s        \\
        3-3t=-14+7s
    \end{cases}$

解得 $\begin{cases}
        t=1 \\
        s=2
    \end{cases}$，求出 $\lambda=2$。

带入，求得交点 $(3,1,0)$

法向量 $n=(2,5,-3)\times(3,-4,7)=\begin{vmatrix}
        i & j  & k  \\
        2 & 5  & -3 \\
        3 & -4 & 7
    \end{vmatrix}=23i-23j-23k=23(1,-1,-1)$

即平面方程为 $(x-3)-(y-1)-z=0$，即 $x-y-z-2=0$
\end{解}

33. 求直线 $\begin{cases} x + y - z - 1 = 0, \\ x - y + z + 1 = 0 \end{cases}$ 在平面 $x + y + z = 0$ 上的投影直线的方程。

\begin{解}
直线 $L_0$ 方向向量 $v_0=(1,1,-1)\times(1,-1,1)=\begin{vmatrix}
        i & j  & k  \\
        1 & 1  & -1 \\
        1 & -1 & 1
    \end{vmatrix}=-2j-2k=-2(0,1,1)$

平面 $A$ 法向量 $n_0=(1,1,1)$

则经过直线 $L_0$ 且垂直于平面 $A$ 得平面 $B$ 得法向量为 $$n_1=(0,1,1)\times(1,1,1)=\begin{vmatrix}
        i & j & k \\
        0 & 1 & 1 \\
        1 & 1 & 1
    \end{vmatrix}=(0,1,-1)$$
$B$ 过 $L_0$ 上一点 $(0,1,0)$，得到平面方程 $B:\,y-z-1=0$

联立 $A,B$ 得所求直线方程 $\begin{cases}
        x+y+z=0 \\
        y-z-1=0
    \end{cases}$
\end{解}

34. 写出垂直于平面 $5x - y + 3z - 2 = 0$，且与它的交线在 $Oxy$ 平面上的平面方程。

\begin{解}
平面 $A$ 与 $Oxy$ 的交线 $L:\,\begin{cases}
        5x - y + 3z - 2 = 0 \\
        z=0
    \end{cases}$

计算其方向向量：$v=(5,-1,3)\times(1,0,0)=\begin{vmatrix}
        i & j  & k \\
        5 & -1 & 3 \\
        0 & 0  & 1
    \end{vmatrix}=-(1,5,0)$

所求平面 $B$ 的法向量：$n=v\times(5,-1,3)=\begin{vmatrix}
        i & j  & k \\
        1 & 5  & 0 \\
        5 & -1 & 3
    \end{vmatrix}=(15,-3,-26)$

$B$ 过 $L$ 上一点 $(1,3,0)$，求出平面方程：$15x-3y-26z-6=0$
\end{解}

\end{document}